\documentclass{homework}
\author{Tashfeen, Ahmad}
\class{CSCI 2114: Data Structures}
\title{Homework 5}
\address{Oklahoma City University, Petree College of Arts \& Sciences, Computer Science}

\graphicspath{{./media}}

\begin{document} \maketitle

\begin{quote}
{\itshape
Est genus in totidem tenui ratione redactum

scriptula, quot menses lubricus annus habet; \\
Parua tabella capit ternos utrimque lapillos,

in qua uicisse est continuasse suos.}
\\\\
There is another game divided into as many parts

as there are months in the year;\\
A small board has three pieces on either side,

the winner must get all the pieces in a straight line.

\hspace{16em}---Ovid (Ars Amatoria III, lines 365-369)
\end{quote}

\question Implement the Minimax algorithm for adversarial search in the \href{https://en.wikipedia.org/wiki/Tic-tac-toe}{Tic-tac-toe} game tree.

Start by \href{https://tashfeen.org/raw/share/ds/minimax.zip}{downloading} and reading the starter code. You need to only edit the \texttt{Minimax.java} by implementing the unfinished Java methods. Upon successfully implementing all of the API, you may run \texttt{javac Game.java} and \texttt{java Game}. You'll be presented to make the first move.  Click at the appropriate place to start the game.

\img<example>[0.14]{Tic-tac-toe gameplay in state of the art graphics by Tashfeen Studios$^\text{TM}$.}{one, two, three, four, five}


\question\label{un} How many unique terminal states are possible for the three by three Tic-tac-toe game?

\question Once your game is functional and your \texttt{Minimax.java} class contains a game tree.  You should start another Java class like this,
\begin{lstlisting}
public class Investigate {
    public static void main(String[] args) {
        Minimax model = new Minimax();
        System.out.println(model.root);
    }
}
\end{lstlisting}

Write appropriate code in this file to answer the following questions,

\begin{enumerate}
    \item How many total leaves does the game tree rooted at \texttt{model.root} have? Is this number greater or less than your answer in question \ref{un}? What does that tell you about the uniqueness of the leaves?
    \item How many tree-leaves result in a draw? A win for the second (min) player? A win for the first (max) player?
\end{enumerate}

\question Ask two other professors to play your implementation on your computer.  At least one of the professors has to be not from computer science.  Ask them if they think they can win and how do they think the ``AI'' maybe picking its moves.  State which professors (and their department) you interviewed and summarise their responses.

\section*{Submission Instructions}

\begin{enumerate}
    \item Turn in a PDF containing any plots, figures and/or answers from the homework.
    \item Turn in your,
    \begin{enumerate}[label=\roman*)]
        \item \texttt{Minimax.java}
        \item \texttt{Investigate.java}
    \end{enumerate}
\end{enumerate}

\end{document}
